% ---
% titre: Lexique - Qui suis-je ? (consignes)
% theme: NO
% chapitre: Nombres naturels et décimaux
% sous_chapitre: 
% niveau: 9VP
% type: JEU
% tags: p-lexique, p-vocabulaire, p-repetition-espacee, f-jeu-regles
% difficulte: 1
% corrige: false
% source: latex
% ---
% NOTE : sous_chapitre laissé vide car ce jeu porte sur le lexique cumulé
% de plusieurs chapitres, pas un sous-chapitre précis. Tags à vérifier
% contre scripts/tags.md.

\header{Lexique --- Qui suis-je\,?}

\textcolor{themeColor}{\textbf{Objectif}}\\
S'entraîner à formuler et reconnaître les définitions du lexique de
mathématiques.

\vspace{20pt}
\textcolor{themeColor}{\textbf{Matériel}}\\
Une feuille et un stylo par élève.

\vspace{20pt}
\textcolor{themeColor}{\textbf{\large Déroulement}}
\begin{enumerate}[leftmargin=*]
\item \textbf{Distribution}
\item[] Chaque élève reçoit \textbf{trois mots} du lexique cumulé.
\medskip
\item \textbf{Rédaction}
\item[]Chaque élève écrit, avec ses propres mots, une définition pour chacun de ses trois termes.
\begin{itemize}[leftmargin=1.2em,label=$\triangleright$]
\item Interdiction de recopier la définition telle quelle depuis le lexique.
\item Interdiction d'utiliser le mot lui-même (ou un mot de la même famille) dans la définition.
\item Un exemple chiffré est autorisé, et même encouragé, s'il aide à cerner le mot sans le nommer.
\end{itemize}
\medskip
\item \textbf{Échange}
\item[]L'élève donne sa feuille de définitions (sans les mots, uniquement les trois définitions numérotées) à un camarade désigné par l'enseignant\,\textperiodcentered\,e.
\item \textbf{Devinette}
\item[] Le camarade lit chaque définition et écrit le mot qu'il pense avoir reconnu.
\item \textbf{Correction croisée}
\item[] Les deux élèves se retrouvent~: l'auteur des définitions annonce les bonnes réponses et explique, si besoin, pourquoi une définition n'a pas permis de deviner le mot (trop vague, mot de la même famille utilisé, etc.).
\end{enumerate}

\vspace{20pt}
\textcolor{themeColor}{\textbf{\large Points}}

\begin{itemize}[leftmargin=*, label=$\diamond$]
\item \textbf{1\,pt} pour chaque mot correctement deviné (pour celui qui devine).
\item \textbf{1\,pt} pour chaque définition qui a permis à l'autre de deviner (pour celui qui a écrit la définition).
\item Un élève peut marquer jusqu'à \textbf{6 points} au total (3 en tant qu'auteur, 3 en tant que devineur).
\end{itemize}
%
%\vspace{6pt}
%\textcolor{themeColor}{\textbf{\large Variante possible}}
%
%Si un mot n'est pas deviné, l'auteur peut donner \textbf{un indice supplémentaire} (mais pas le mot) avant une deuxième tentative --- au
%prix d'une pénalité d'un demi-point, pour continuer à travailler le mot
%sans le braquer sur l'échec.